\section*{Capítulo \ref{ch:b2:reaction-enthalpy} --- Entalpias de reação}

\begin{solution}{exo:b2:reaction-enthalpy:1}
$\Delta_r H^\circ = 3(-393.51) + 4(-285.83) - (-104.7) - 5(0) =
\qty{-2219.2}{kJ/mol}$, negativa: \omterm{def:b2:reaction-enthalpy:exothermic}{exotérmica}. (A entalpia de combustão
medida é a mesma, \qty{-2219.2}{kJ/mol}.)
\end{solution}

\begin{solution}{exo:b2:reaction-enthalpy:2}
Metano: $890.3/16.04 = \qty{55.5}{kJ/g}$, ou seja, \qty{55.5}{MJ} por
quilograma. Di-hidrogênio: $285.83/2.016 = \qty{141.8}{kJ/g}$,
\qty{141.8}{MJ/kg}: cerca de 2.6 vezes mais por quilograma (mas muito menos
por litro, já que o gás é leve).
\end{solution}

\begin{solution}{exo:b2:reaction-enthalpy:3}
$\Delta\nu_{\text{gás}} = 6 - 7.5 = -1.5$, logo $\Delta_r H = \Delta_r U +
\Delta\nu_{\text{gás}}RT = -3264 - 1.5 \times 8.314 \times 298.15 \times
10^{-3} = \qty{-3267.7}{kJ/mol}$, de acordo com o valor calculado a partir
das entalpias de formação, \qty{-3267.6}{kJ/mol}.
\end{solution}

\begin{solution}{exo:b2:reaction-enthalpy:4}
A reação-alvo é a primeira reação menos duas vezes a segunda:
$\Delta_r H^\circ = -393.5 - 2(-283.0) = \qty{+172.5}{kJ/mol}$. Ela é
\omterm{def:b2:reaction-enthalpy:exothermic}{endotérmica}: o coque quente esfria ao transformar o dióxido de carbono em
monóxido de carbono.
\end{solution}

\begin{solution}{exo:b2:reaction-enthalpy:5}
Ligações rompidas: C--H e Cl--Cl, $416 + 243 = \qty{659}{kJ/mol}$; ligações
formadas: C--Cl e H--Cl, $350 + 432 = \qty{782}{kJ/mol}$: estimativa
\qty{-123}{kJ/mol}. A partir das entalpias de formação: $-83.7 - 92.3 + 74.9 =
\qty{-101.1}{kJ/mol}$. A diferença de \qty{22}{kJ/mol} vem do valor médio
de C--Cl, uma média sobre várias moléculas, e da ligação C--H rompida no
metano, mais forte que a média das quatro.
\end{solution}

\begin{solution}{exo:b2:reaction-enthalpy:6}
A \qty{298}{K}: $-241.83 + 285.83 = \qty{44.00}{kJ/mol}$. Kirchhoff com
$\Delta C_p = 33.59 - 75.35 = \qty{-41.76}{J/(K.mol)}$:
$44.00 - 0.04176 \times 101.85 = \qty{39.75}{kJ/mol}$ a \qty{400}{K}. As
tabelas dão $-242.85 + 282.59 = \qty{39.74}{kJ/mol}$: as capacidades
caloríficas constantes são excelentes nesse intervalo.
\end{solution}

\begin{solution}{exo:b2:reaction-enthalpy:7}
Ionização: $4.341 \times 96.485 = \qty{418.8}{kJ/mol}$; captura:
$-3.613 \times 96.485 = \qty{-348.6}{kJ/mol}$. \omterm{def:b2:reaction-enthalpy:lattice-enthalpy}{Entalpia de rede} $= 436.7 +
89.0 + 121.3 + 418.8 - 348.6 = \qty{717}{kJ/mol}$, menor que os
\qty{787}{kJ/mol} do cloreto de sódio: \ce{K+} é maior que \ce{Na+}, os
íons ficam mais afastados e se atraem menos.
\end{solution}

\begin{solution}{exo:b2:reaction-enthalpy:8}
Capacidade calorífica do calorímetro: $C = 2.00/0.418 = \qty{4.785}{kJ/K}$.
Calor liberado: $4.785 \times 0.583 = \qty{2.789}{kJ}$, para $\xi =
\qty{0.0500}{mol}$ (o ácido é o limitante, a base está em leve excesso):
$\Delta_r H = -2.789/0.0500 = \qty{-55.8}{kJ/mol}$.
\end{solution}

\begin{solution}{exo:b2:reaction-enthalpy:9}
\ce{C8H18(l) + 25/2O2(g) -> 8CO2(g) + 9H2O(l)}: $\Delta_r H^\circ =
8(-393.51) + 9(-285.83) + 250.3 = \qty{-5470.3}{kJ/mol}$; com $M =
\qty{114.2}{g/mol}$, \qty{47.9}{MJ/kg}. Dióxido de carbono: $8 \times 44.01 =
\qty{352.1}{g}$ para \qty{5.470}{MJ}, isto é, \qty{64.4}{g/MJ}, contra
$44.01/0.8903 = \qty{49.4}{g/MJ}$ para o metano: o metano tem mais hidrogênio
por carbono, e queimar hidrogênio libera calor sem dióxido de carbono.
\end{solution}

\begin{solution}{exo:b2:reaction-enthalpy:10}
Produtos de \ce{CH4 + 2O2 -> CO2 + 2H2O(g)}: $C_p = 37.13 + 2 \times 33.59 =
\qty{104.3}{J/K}$; $q = \qty{802.3}{kJ}$; $T_f = 298 + 802\,300/104.3 \approx
\qty{7990}{K}$, quase três vezes o valor no ar, porque nenhum nitrogênio
fica com sua parte do calor. A chama real no oxigênio é muito mais fria:
muito antes de \qty{7990}{K}, o dióxido de carbono e a água
se dissociam em monóxido de carbono, hidrogênio, oxigênio e átomos, reações
\omterm{def:b2:reaction-enthalpy:exothermic}{endotérmicas} que absorvem a maior parte do calor; e as capacidades
caloríficas crescem com $T$.
\end{solution}

\begin{solution}{exo:b2:reaction-enthalpy:11}
$\Delta_r H^\circ(700) = \Delta_r H^\circ(298) + \int_{298}^{700}(\Delta_r a +
\Delta_r b\,T)\,\dd T = -91.88 + [-60.5 \times 401.85 + 0.0270 \times
(700^2 - 298.15^2)] \times 10^{-3} = -91.88 - 24.31 + 10.83 =
\qty{-105.4}{kJ/mol}$, a \qty{0.2}{kJ/mol} das tabelas
(\qty{-105.2}{kJ/mol}), enquanto a estimativa com $C_p$ constante errava por
\qty{4.5}{kJ/mol}.
\end{solution}

\begin{solution}{exo:b2:reaction-enthalpy:12}
$\Delta_r H^\circ = 2 \times 90.29 = \qty{+180.6}{kJ/mol}$ (\omterm{def:b2:reaction-enthalpy:exothermic}{endotérmica}).
$\Delta_r C_p^\circ = 2(29.85) - 29.12 - 29.38 = \qty{1.2}{J/(K.mol)}$; logo,
entre \qty{298}{K} e \qty{2000}{K}, a \omterm{def:b2:reaction-enthalpy:reaction-quantity}{entalpia de reação} varia de cerca de
$1.2 \times 1702 = \qty{2}{kJ/mol}$, um por cento. Uma reação \omterm{def:b2:reaction-enthalpy:exothermic}{endotérmica}
é favorecida pela alta temperatura (\cref{ch:b2:equilibrium-shifts}): o ar
frio quase não forma monóxido de nitrogênio, as chamas mais quentes formam
mais; por isso a temperatura de combustão é uma alavanca contra esse
poluente.
\end{solution}

\begin{solution}{pb:b2:reaction-enthalpy:1}
\textbf{1.} \ce{C4H10(g) + 13/2O2(g) -> 4CO2(g) + 5H2O(l)}.
\textbf{2.} $\Delta_r H^\circ = 4(-393.51) + 5(-285.83) + 125.6 =
\qty{-2877.6}{kJ/mol}$.
\textbf{3.} Com vapor d'água, $4(-393.51) + 5(-241.83) + 125.6 =
\qty{-2657.6}{kJ/mol}$; a diferença, \qty{220.0}{kJ}, é a entalpia de
vaporização de cinco mols de água (\qty{44.0}{kJ/mol} cada).
\textbf{4.} $2877.6/58.12 = \qty{49.5}{kJ/g}$ e $2657.6/58.12 =
\qty{45.7}{kJ/g}$.
\textbf{5.} O segundo: a água sai como vapor e seu calor de condensação é
perdido.
\textbf{6.} $230 \times 45.7 = \qty{1.05e4}{kJ}$, cerca de \qty{10.5}{MJ}.
\textbf{7.} A volume constante, o calor recebido é $\Delta U$ (não há
trabalho de pressão): a bomba mede $\Delta_r U$.
\textbf{8.} $\Delta\nu_{\text{gás}} = 4 - 1 - 6.5 = -3.5$ (a água é
líquida).
\textbf{9.} $\Delta_r U = \Delta_r H - \Delta\nu_{\text{gás}}RT = -2877.6 +
3.5 \times 2.479 = \qty{-2868.9}{kJ/mol}$.
\textbf{10.} $\xi = 0.500/58.12 = \qty{8.60e-3}{mol}$; calor $8.60 \times
10^{-3} \times 2868.9 = \qty{24.68}{kJ}$; $\Delta T = 24.68/10.00 =
\qty{2.468}{K}$.
\textbf{11.} Calor $10.00 \times 2.461 = \qty{24.61}{kJ}$: $\Delta_r U =
-24.61/(8.603 \times 10^{-3}) = \qty{-2860.6}{kJ/mol}$ e $\Delta_r H =
-2860.6 - 8.7 = \qty{-2869.3}{kJ/mol}$, \qty{0.29}{\percent} abaixo das
tabelas em valor absoluto: calor perdido pelo calorímetro ou uma combustão
ligeiramente incompleta.
\textbf{12.} $n = 1000/18.02 = \qty{55.5}{mol}$; $Q = 55.5 \times 75.35
\times 80 = \qty{334.5}{kJ}$.
\textbf{13.} Calor liberado $16.0 \times 45.72 = \qty{731.5}{kJ}$; rendimento
$334.5/731.5 = 0.46$.
\textbf{14.} Para os gases de combustão quentes que escapam em volta da
panela, o ar aquecido pela chama e as paredes da panela, e a radiação.
\textbf{15.} $230 \times 45.72 \times 0.457/334.5 = \qty{14.4}{L}$.
\textbf{16.} \qty{6.5}{mol} de \ce{O2} vêm com $6.5 \times 78.08/20.95 =
\qty{24.23}{mol}$ de \ce{N2} e $6.5 \times 0.93/20.95 = \qty{0.289}{mol}$
de \ce{Ar}.
\textbf{17.} \qty{4}{mol} \ce{CO2}, \qty{5}{mol} \ce{H2O(g)},
\qty{24.23}{mol} \ce{N2}, \qty{0.289}{mol} \ce{Ar}.
\textbf{18.} A chama é adiabática e isobárica, logo $\Delta H = 0$. Ao longo
de um caminho formado pela reação a \qty{298}{K} ($\Delta H_1 =
\qty{-2657.6}{kJ}$) seguida do aquecimento desses produtos, $\Delta H_2 =
-\Delta H_1 = \qty{+2657.6}{kJ}$: todo o calor vai para os gases.
\textbf{19.} $C_p = 4(37.13) + 5(33.59) + 24.23(29.12) + 0.289(20.79) =
\qty{1028}{J/K}$.
\textbf{20.} $T_f = 298 + 2\,657\,600/1028 = \qty{2883}{K}$.
\textbf{21.} $T_f = 2300 + 100 \times (2657.6 - 2515.7)/(2655.7 - 2515.7) =
\qty{2401}{K}$.
\textbf{22.} As capacidades caloríficas dos gases crescem com a temperatura
(de cerca de um quarto entre \qty{298}{K} e \qty{2400}{K}): os valores à
temperatura ambiente aquecem os gases barato demais.
\textbf{23.} Acima de \qty{2000}{K}, o dióxido de carbono e a água se
dissociam parcialmente (\omterm{def:b2:reaction-enthalpy:exothermic}{endotérmico}), a chama irradia, e a mistura com o ar
nunca é perfeita.
\textbf{24.} A \omterm{def:b2:reaction-enthalpy:flame-temperature}{temperatura adiabática de chama} do butano no ar é
$\boldsymbol{T_{\text{ad}} \approx \qty{2.40e3}{K}}$.
\end{solution}
